\scanpage{33}
\mspage{32}
[Il s'ensuit qu'on a \unsure{dans|donc} $\unsure{M|\mathcal M}\simeq M^{*}$ une relation d'ordre $\leq$, avec
\[
\begin{aligned}
 X\leq Y\quad\Longleftrightarrow\quad&
 \deleted{\mathrm{Omb}(x)}\Phi_{\mathrm{Omb}(x)}
 \subset\Phi_{\mathrm{Omb}(Y)}\\
 &\qquad\deleted{\text{i.e.}}\ \text{i.e.}\unsure{,|.}\quad
 \mathrm{Omb}(x)\in
 \underset{\inserted{\substack{\mathrm{def}\ \unsure{\shortparallel|\text{''}}\\
 \unsure{\mathrm{Mulomb}|\mathrm{Multomb}}(Y)}}}{\Phi_{\mathrm{Omb}(Y)}},
\end{aligned}
\]
de telle façon que les \deleted{éléments de} $F\in\widetilde{\unsure{\mathcal F|\mathfrak F|\mathscr F}}$ soient des parties fermées de $M$\unsure{,|.}
Associant à tout $X\in M$ la partie $\deleted{\unsure{\widetilde X|X}}$
\[
\begin{aligned}
 \widetilde X&=\{Y\in M\mid Y\leq\unsure{x|X}\ \text{i.e.}\quad
 \mathrm{Omb}(Y)\in\Phi_{\mathrm{Omb}(x)}\}\\
 &=M_{\leq x}
\end{aligned}
\]
on trouve une application injective
\[
 \overset{\unsure{\cdot\,\cdot}}{M}\hookrightarrow\widetilde{\mathcal F}\quad(\simeq\mathcal F_M)
\]
qui induit un iso entre $M$ et l'ens des éléments str. irréductibles de $\widetilde{\mathcal F}$ (ordonné par inclusion, notée $\leq$).]

\deleted{\unsure{Donc|D}}

\uline{Théorème} - \unsure{\uline{Scholie}|\uuline{Scholie}}\quad La catégorie \inserted{(groupoïde)} des \deleted{\unsure{$\mathcal F$}} ateliers $(\mathcal F,\leq,\ll)$ satisfaisant à \deleted{\unsure{C}}At 1 \unsure{.|à|,} At 7, équivaut à celle des $(M,\mathcal F_M)$ (avec $\mathcal F_M\in\deleted{\unsure{\mathcal P|\mathfrak P|\mathscr P|P}\mathcal P}\mathrm{Fig}(M)$) satisfaisant les conditions suivantes

\uline{At$^{*}$1}, $\mathcal F_M$ \deleted{non vide, et} \inserted{contient la figure vide, et} stable par \deleted{\unsure{les sup majorés|sup majorés}} \inserted{réunions} majorées (dans $\mathcal F_M$)

\deleted{\unsure{Cela|Ceci}} [La stabilité par sup \unsure{équivaudra|équivaudre}, moyennant la condition suivante, à celle de \deleted{\unsure{passage}} stabilité par passage à une \uline{sous} \deleted{\unsure{famille fermée|famille finie}} \uline{figure}.)
